% =====================================================================
% PATCH: replaces \subsection{Sampled policy fitting and dual update}
%        and the algorithm environment alg:nbpo in main_v5.tex.
%
% Why: the alternating fixed-opponent map that the current text presents as
% THE practical inner solve is non-convergent in exactly the regime the method
% is for. Measured on the feasibility-preserving controlled benchmark, it
% converged on 0 of 25 cells and reached a normalized min surplus of -3.373 at
% the highest intransitivity, against +0.931 for the direct solve. It therefore
% moves to the appendix as an approximation and failure ablation, and the
% direct prompt-separable concave solve becomes the reported practical method.
%
% What is NOT changed: the population theorem, and every statement about the
% exact population proximal update.
% =====================================================================

\subsection{Finite-pool inner solve and dual update}
\label{sec:finite-pool-inner}
Fix outer stage $t$ and a dual vector $\lambda$. On the frozen finite response
pool, \eqref{eq:weighted-policy} has structure that makes it solvable rather
than merely approximable. Writing $V_{k,\beta_k}(\pi)=\E_x[v_{k,x}(\pi(\cdot\mid x))]$,
each $v_{k,x}$ is a soft-min of terms linear in $\pi(\cdot\mid x)$ and therefore
\emph{concave}, and it depends on $\pi$ only through that prompt's row. The
stage objective
\[
J_\lambda(\pi)
=
\sum_{k=1}^K \lambda_k s_k(\pi)-\tfrac{1}{\eta_t}D(\pi\,\|\,\pi_t)
\]
is thus concave and \emph{separates completely across prompts}: it is $|\mathcal{X}|$
independent concave programs over the response simplex, not one program over
their product. We solve them directly, in double precision, to a declared
stationarity tolerance of $10^{-4}$.

We return the update \eqref{eq:fixed-opponent-update} evaluated \emph{at} that
maximizer rather than the maximizer itself. This is not cosmetic: it makes the
log-ratio identity that \eqref{eq:binary-target} depends on hold to machine
precision by construction (measured residual $\le 1.3\times10^{-15}$), whereas an
optimizer iterate satisfies it only to its own tolerance.

The dual is likewise solved rather than descended. Stationarity of
\eqref{eq:prox-dual} is $s_k(\pi(\lambda))=1/\lambda_k$, and $\pi(\lambda)$ is
a well-defined map once the inner problem is solved, so this is a $K$-dimensional
root problem, which we solve in $\log\lambda$ coordinates and stop on the
projected KKT residual at $10^{-6}$ rather than after a fixed iteration count.

Both choices matter numerically. The raw multipliers satisfy $\lambda_k=1/s_k$
and therefore diverge as a surplus approaches zero, which is precisely the regime
a bargaining solution occupies; a projected subgradient stalls near a residual of
$10^{-2}$ there, and the alternating inner map ceases to contract entirely
(Appendix~\ref{app:rstep-ablation}).

\paragraph{Cost.} The separability is what makes this practical at scale: the
per-prompt programs are independent, so the solve is linear in the number of
prompts and embarrassingly parallel. At $7000$ prompts with a $4{+}4$ pool, one
complete dual solve takes $10.6$\,s over $28$ dual evaluations at a peak resident
size of $534$\,MB; at $16{+}16$ it takes $15.3$\,s. All $24$ measured
configurations converged, with projected KKT residuals between $1.8\times10^{-15}$
and $7.3\times10^{-10}$.

\paragraph{What is and is not exact.} Three levels should not be conflated. The
population proximal update of \eqref{eq:prox-subproblem} is exact by
Theorem~\ref{thm:last-iterate}. The finite-pool inner problem is \emph{solved
numerically to a declared tolerance}. The projection of the finite-pool
log-ratio target into the neural policy by pairwise regression remains
\emph{approximate}, and is the approximation the method carries. We therefore
describe the middle step as a \emph{direct finite-pool concave inner solve}, and
never describe the algorithm as a whole as exact.

% ---------------------------------------------------------------------
\begin{algorithm}[t]
\caption{Nash Bargaining Preference Optimization}
\label{alg:nbpo}
\begin{algorithmic}[1]
\REQUIRE Reference $\piref$, preference models $P_{1:K}$, temperatures $\beta_{1:K}$, outer stages $T$, step $\eta_t$, dual tolerance $\epsilon_{\mathrm{dual}}$, inner tolerance $\epsilon_{\mathrm{in}}$.
\STATE Estimate $d_k=V_{k,\beta_k}(\piref)$ and set $\pi_0\leftarrow\piref$.
\FOR{$t=0,\ldots,T-1$}
    \STATE Sample the response pool at $\pi_t$ and build the centered tensors $A_k$. \hfill\COMMENT{Eq.~\eqref{eq:centered}}
    \STATE Initialize $\lambda^{(0)}>0$ (warm-start from the previous stage when available).
    \STATE \textbf{Solve} $s_k(\pi(\lambda))=1/\lambda_k$ for $\lambda^\star$, stopping at projected KKT residual $\le\epsilon_{\mathrm{dual}}$, where for each trial $\lambda$:
    \STATE \quad for every prompt $x$ \textbf{in parallel}: $\;\widehat\pi_x\leftarrow\arg\max_{p\in\Delta}\;\eta_t\sum_k\lambda_k v_{k,x}(p)-\mathrm{KL}(p\,\|\,\pi_t(\cdot\mid x))$, to $\epsilon_{\mathrm{in}}$
    \STATE \quad $\pi(\lambda)\leftarrow$ \eqref{eq:fixed-opponent-update} evaluated at $\widehat\pi$; \; $s\leftarrow V(\pi(\lambda))-d$
    \STATE Form the finite-pool target $\log\pi^\star-\log\pi_t$ from $\lambda^\star$ and its opponent.
    \STATE $\widetilde\pi_{t+1}\leftarrow$ one pairwise regression fit of \eqref{eq:regression-loss} onto that target. \hfill\COMMENT{the approximate step}
    \STATE Set $\pi_{t+1}\leftarrow\widetilde\pi_{t+1}$ if all held-out empirical surpluses are positive; otherwise retain $\pi_t$.
\ENDFOR
\STATE \textbf{return} $\pi_T$.
\end{algorithmic}
\end{algorithm}
% ---------------------------------------------------------------------
The held-out check is an implementation safeguard rather than a finite-sample
guarantee. The convergence result below concerns exact population solutions of
\eqref{eq:prox-subproblem} and is unaffected by how the finite-pool inner problem
is solved. Exactly one neural fit is performed per outer stage, after the dual
has converged; the dual loop itself touches no network.
